\documentclass[CJK,13pt]{beamer}
\usepackage{CJKutf8}
\usepackage{beamerthemesplit}
\usetheme{Malmoe}
\useoutertheme[footline=authortitle]{miniframes}
\usepackage{amsmath}
\usepackage{amssymb}
\usepackage{graphicx}
\usepackage{eufrak}
\usepackage{color}
\usepackage{slashed}
\usepackage{simplewick}
\usepackage{tikz}
\usepackage{tcolorbox}
\usepackage{ulem}
\graphicspath{{../figures/}}
%%figures
\def\lfig#1#2{\includegraphics[width=#1 in]{#2}}
\def\tfig#1#2{\includegraphics[height=#1 in]{#2}}
\def\addfig#1#2{\begin{center}\includegraphics[width=#1 in]{#2}\end{center}}
\def\question{\includegraphics[width=0.3in]{why.jpg}\,}
\def\answer{\includegraphics[width=0.3in]{baozou_haha.png}\,}
\def\wulian{\includegraphics[width=0.18in]{emoji_wulian.jpg}}
\def\bigwulian{\includegraphics[width=0.35in]{emoji_wulian.jpg}}
\def\bye{\includegraphics[width=0.18in]{emoji_bye.jpg}}
\def\bigbye{\includegraphics[width=0.35in]{emoji_bye.jpg}}
\def\huaixiao{\includegraphics[width=0.18in]{emoji_huaixiao.jpg}}
\def\bighuaixiao{\includegraphics[width=0.35in]{emoji_huaixiao.jpg}}
\def\jianxiao{\includegraphics[width=0.18in]{emoji_jianxiao.jpg}}
\def\bigjianxiao{\includegraphics[width=0.35in]{emoji_jianxiao.jpg}}
\def\haoqi{\includegraphics[width=0.18in]{emoji_haoqi.jpg}}
%% colors
\def\blacktext#1{{\color{black}#1}}
\def\bluetext#1{{\color{blue}#1}}
\def\redtext#1{{\color{red}#1}}
\def\darkbluetext#1{{\color[rgb]{0,0.2,0.6}#1}}
\def\skybluetext#1{{\color[rgb]{0.2,0.7,1.}#1}}
\def\cyantext#1{{\color[rgb]{0.,0.5,0.5}#1}}
\def\greentext#1{{\color[rgb]{0,0.7,0.1}#1}}
\def\darkgray{\color[rgb]{0.2,0.2,0.2}}
\def\lightgray{\color[rgb]{0.6,0.6,0.6}}
\def\gray{\color[rgb]{0.4,0.4,0.4}}
\def\blue{\color{blue}}
\def\red{\color{red}}
\def\orange{\color[rgb]{1.,0.8,0.}}
\def\green{\color{green}}
\def\darkgreen{\color[rgb]{0,0.4,0.1}}
\def\darkblue{\color[rgb]{0,0.2,0.6}}
\def\skyblue{\color[rgb]{0.2,0.7,1.}}
%%control
\def\diag{\mathrm{diag}\,}
\def\heaviside{\,\mathrm{h}\,}
\def\bral{\left(\begin{array}{l}}
\def\brar{\end{array}\right)}
\def\brall{\left(\begin{array}{ll}}
\def\brarr{\end{array}\right)}
\def\bralll{\left(\begin{array}{lll}}
\def\brarrr{\end{array}\right)}
\def\branchl{\left\{\begin{array}{l}}
\def\branchr{\end{array}\right.}
\def\branchll{\left\{\begin{array}{ll}}
\def\branchrr{\end{array}\right.}
\def\branchlll{\left\{\begin{array}{lll}}
\def\branchrrr{\end{array}\right.}
\def\sfgamma#1{\,\Gamma\left( #1 \right)\,}
\def\be{\begin{equation}}
\def\ee{\nonumber\end{equation}}
\def\bea{\begin{eqnarray}}
\def\eea{\nonumber\end{eqnarray}}
\def\bch{\begin{CJK}{UTF8}{gbsn}}
\def\ech{\end{CJK}}
\def\bitem{\begin{itemize}}
\def\eitem{\end{itemize}}
\def\bcenter{\begin{center}}
\def\ecenter{\end{center}}
\def\bex{\begin{minipage}{0.2\textwidth}\includegraphics[width=0.6in]{jugelizi.png}\end{minipage}\begin{minipage}{0.76\textwidth}}
\def\eex{\end{minipage}}
\def\chtitle#1{\frametitle{\bch#1\ech}}
\def\bmat#1{\left(\begin{array}{#1}}
\def\emat{\end{array}\right)}
\def\bcase#1{\left\{\begin{array}{#1}}
\def\ecase{\end{array}\right.}
\def\bmini#1{\begin{minipage}{#1\textwidth}}
\def\emini{\end{minipage}}
\def\tbox#1{\begin{tcolorbox}#1\end{tcolorbox}}
\def\pfrac#1#2#3{\left(\frac{\partial #1}{\partial #2}\right)_{#3}}
\def\res#1#2{\,\mathrm{res}\,#1\left(#2\right)\,}
\def\newt#1#2{\left(\begin{array}{c}#1\\ #2\end{array}\right)}
\def\reof#1{\,\mathrm{Re}{\left(#1\right)\,}}
\def\imof#1{\,\mathrm{Im}{\left(#1\right)\,}}
\def\Arg#1{\,\mathrm{Arg}\,#1\,}
%%symbols
\def\intfull{\int_{-\infty}^\infty}
\def\inthalf{\int_0^\infty}
\def\bropt{\,(\ \ \ )}
\def\sech{\mathrm{sech}\,}
\def\csch{\mathrm{csch}\,}
\def\asinh{\mathrm{asinh}\,}
\def\acosh{\mathrm{acosh}\,}
\def\sone{$\star$}
\def\stwo{$\star\star$}
\def\sthree{$\star\star\star$}
\def\sfour{$\star\star\star\star$}
\def\sfive{$\star\star\star\star\star$}
\def\rint{{\int_\leftrightarrow}}
\def\roint{{\oint_\leftrightarrow}}
\def\stdHf{{\textit{\r H}_f}}
\def\deltaH{{\Delta \textit{\r H}}}
\def\ii{{\dot{\imath}}}
\def\skipline{{\vskip0.1in}}
\def\skiplines{{\vskip0.2in}}
\def\lagr{{\mathcal{L}}}
\def\hamil{{\mathcal{H}}}
\def\vecT{{\mathbf{T}}}
\def\vecN{{\mathbf{N}}}
\def\vecB{{\mathbf{B}}}
\def\vecv{{\mathbf{v}}}
\def\vecr{{\mathbf{r}}}
\def\vecf{{\mathbf{f}}}
\def\vecg{{\mathbf{g}}}
\def\vecupsilon{{\mathbf{\upsilon}}}
\def\vecu{{\mathbf{u}}}
\def\vecj{{\mathbf{j}}}
\def\vecx{{\mathbf{x}}}
\def\vecy{{\mathbf{y}}}
\def\vecz{{\mathbf{z}}}
\def\veck{{\mathbf{k}}}
\def\vecp{{\mathbf{p}}}
\def\vecn{{\mathbf{n}}}
\def\vecA{{\mathbf{A}}}
\def\vecP{{\mathbf{P}}}
\def\vecsigma{{\mathbf{\sigma}}}
\def\hatJn{{\hat{J_\vecn}}}
\def\hatJx{{\hat{J_x}}}
\def\hatJy{{\hat{J_y}}}
\def\hatJz{{\hat{J_z}}}
\def\hatj#1{\hat{J_{#1}}}
\def\hatphi{{\hat{\phi}}}
\def\hatq{{\hat{q}}}
\def\hatpi{{\hat{\pi}}}
\def\vel{\upsilon}
\def\Dint{{\mathcal{D}}}
\def\adag{{\hat{a}^\dagger}}
\def\bdag{{\hat{b}^\dagger}}
\def\cdag{{\hat{c}^\dagger}}
\def\ddag{{\hat{d}^\dagger}}
\def\hata{{\hat{a}}}
\def\hatb{{\hat{b}}}
\def\hatc{{\hat{c}}}
\def\hatd{{\hat{d}}}
\def\hatD{{\,\hat{D}}}
\def\hatN{{\hat{N}}}
\def\hatH{{\hat{H}}}
\def\hatp{{\hat{p}}}
\def\Fup{{F^{\mu\nu}}}
\def\Fdown{{F_{\mu\nu}}}
\def\newl{\nonumber \\}
\def\vece{\mathrm{e}}
\def\calM{{\mathcal{M}}}
\def\calT{{\mathcal{T}}}
\def\calR{{\mathcal{R}}}
\def\barpsi{\bar{\psi}}
\def\baru{\bar{u}}
\def\barv{\bar{\upsilon}}
\def\qeq{\stackrel{?}{=}}
\def\ftf{\stackrel{\mathcal{FT}}{\Longrightarrow}}
\def\ftb{\stackrel{\mathcal{FT}}{\Longleftarrow}}
\def\ftfb{\stackrel{\mathcal{FT}}{\Longleftrightarrow}}
\def\ltf{\stackrel{\mathcal{LT}}{\Longrightarrow}}
\def\ltb{\stackrel{\mathcal{LT}}{\Longleftarrow}}
\def\ltfb{\stackrel{\mathcal{LT}}{\Longleftrightarrow}}
\def\torder#1{\mathcal{T}\left(#1\right)}
\def\rorder#1{\mathcal{R}\left(#1\right)}
\def\contr#1#2{\contraction{}{#1}{}{#2}#1#2}
\def\trof#1{\mathrm{Tr}\left(#1\right)}
\def\trace{\mathrm{Tr}}
\def\comm#1{\ \ \ \left(\mathrm{used}\ #1\right)}
\def\tcomm#1{\ \ \ (\text{#1})}
\def\slp{\slashed{p}}
\def\slk{\slashed{k}}
\def\calp{{\mathfrak{p}}}
\def\veccalp{\mathbf{\mathfrak{p}}}
\def\Tthree{T_{\tiny \textcircled{3}}}
\def\pthree{p_{\tiny \textcircled{3}}}
\def\dbar{{\,\mathchar'26\mkern-12mu d}}
\def\erf{\mathrm{erf}}
\def\const{\mathrm{const.}}
\def\pheat{\pfrac p{\ln T}V}
\def\vheat{\pfrac V{\ln T}p}
%%units
\def\fdeg{{^\circ \mathrm{F}}}
\def\cdeg{^\circ \mathrm{C}}
\def\atm{\,\mathrm{atm}}
\def\angstrom{\,\text{\AA}}
\def\SIL{\,\mathrm{L}}
\def\SIkm{\,\mathrm{km}}
\def\SIyr{\,\mathrm{yr}}
\def\SIGyr{\,\mathrm{Gyr}}
\def\SIV{\,\mathrm{V}}
\def\SImV{\,\mathrm{mV}}
\def\SIeV{\,\mathrm{eV}}
\def\SIkeV{\,\mathrm{keV}}
\def\SIMeV{\,\mathrm{MeV}}
\def\SIGeV{\,\mathrm{GeV}}
\def\SIcal{\,\mathrm{cal}}
\def\SIkcal{\,\mathrm{kcal}}
\def\SImol{\,\mathrm{mol}}
\def\SIN{\,\mathrm{N}}
\def\SIHz{\,\mathrm{Hz}}
\def\SIm{\,\mathrm{m}}
\def\SIcm{\,\mathrm{cm}}
\def\SIfm{\,\mathrm{fm}}
\def\SImm{\,\mathrm{mm}}
\def\SInm{\,\mathrm{nm}}
\def\SImum{\,\mathrm{\mu m}}
\def\SIJ{\,\mathrm{J}}
\def\SIW{\,\mathrm{W}}
\def\SIkJ{\,\mathrm{kJ}}
\def\SIs{\,\mathrm{s}}
\def\SIkg{\,\mathrm{kg}}
\def\SIg{\,\mathrm{g}}
\def\SIK{\,\mathrm{K}}
\def\SImmHg{\,\mathrm{mmHg}}
\def\SIPa{\,\mathrm{Pa}}
%page
\def\secpage#1#2{\begin{frame}\bcenter{\bf \Huge #1} \ecenter \skipline \tbox{\bcenter #2 \ecenter}\end{frame}}
\def\append#1#2{\secpage{附录#1}{#2}}
\def\thinka#1{\begin{frame}\frametitle{思考题}\bcenter\lfig{0.4}{think0.jpg}  \skipline #1 \ecenter\end{frame}}
\def\thinkb#1{\begin{frame}\frametitle{思考题}\bcenter\lfig{0.5}{think1.jpg}  \skipline #1 \ecenter\end{frame}}
\def\thinkc#1{\begin{frame}\frametitle{思考题}\bcenter\lfig{0.5}{think2.jpg}  \skipline #1 \ecenter\end{frame}}
\def\thinkd#1{\begin{frame}\frametitle{思考题}\bcenter\lfig{0.5}{think3.jpg}  \skipline #1 \ecenter\end{frame}}
\def\thinke#1{\begin{frame}\frametitle{思考题}\bcenter\lfig{0.5}{think4.jpg}  \skipline #1 \ecenter\end{frame}}
\def\thinkf#1{\begin{frame}\frametitle{思考题}\bcenter\lfig{0.5}{think5.jpg}  \skipline #1 \ecenter\end{frame}}
\def\schw{ds^2 = \left(1-\frac{2GM}{r}\right)dt^2 - \left(1-\frac{2GM}{r}\right)^{-1} dr^2 - r^2\left(d\theta^2 + \sin^2\theta d\phi^2\right)}

\input{titlepage.tex}
  \date{}
  \begin{document}
  \bch
\tpage{9}{Four-Dimensional Spacetime}


\begin{frame}
  \frametitle{从二维曲面到四维时空}

  二维曲面的内蕴几何学已经具有足够的代表性。在大多数情况下阿，可以直接把二维曲面的内蕴几何学的结论推广到四维时空。但是，我们之前很多的推导都借助了嵌入高维平直空间的假设，下面我们再以尽量不借助高维平直空间的形式过一遍这些概念，以加深对广义相对论的逻辑架构的理解。

  \skiplines

  为了和大多数文献保持一致，我们将使用 $(x^0, x^1, x^2, x^3)$ 来标记四维时空的坐标。这里的 $x^0$ 经常具有类似于“时间”的意义——但要注意弯曲时空的时间和空间坐标选取具有更大的随意性，一般不能直接用狭义相对论那套“默认上帝视角”的术语来描述问题。

  {\scriptsize 你或者清晰地说“xx事件在xx坐标系的时间坐标”，或者明确地描述事件“xx通过望远镜观测到xx钟的时间”，模棱两可的描述“对xx观测者来说xx的钟走得更慢”在广义相对论里是不被认可的。}

\end{frame}

\secpage{度规}{$$ds^2=g_{\mu\nu}dx^\mu dx^\nu$$}

\begin{frame}
  \frametitle{Minkowski时空}
  狭义相对论里的 Minkowski 时空是四维时空的一种特殊情况。如果选取惯性运动者的固连直角坐标系，Minkowski 时空的第一基本形式——以后我们简单就说它的度规是：
  \be
  ds^2 = (dx^0)^2 - (dx^1)^2 - (dx^2)^2 - (dx^3)^2.
  \ee
  也就是度规矩阵是个对角矩阵 $\diag(1,-1,-1,-1)$。

  \skipline
  
  {\scriptsize 在很多文献中，也有把 Minkowski 时空的度规取为 $\diag(-1, 1, 1, 1)$ 的。两种取法并无谁对谁错之分，但是在同一个工作里，你最好始终坚持用一种取法，否则计算中很容易产生符号错误。}
  
\end{frame}


\begin{frame}
  \frametitle{Minkowski度规}
  {\blue 由对角矩阵 $\diag(1,-1,-1,-1)$ 定义的 Minkowski 度规写作 $\eta_{\mu\nu}$。}
  这样在 Minkowski 时空 $ ds^2 = \eta_{\mu\nu} dx^\mu dx^\nu.$
  
  \addfig{1.9}{yueyueyue.jpg}
  
  从这一讲开始，我们将约定：{\blue 默认用希腊字母来标记四维时空指标，用拉丁字母来标记三维空间指标。} 例如 $p^\mu p_\mu$ 默认表示 $p^0p_0+p^1p_1+p^2p_2+p^3p_3$，而 $p^ip_i$ 默认表示 $p^1p_1 + p^2p_2 + p^3p_3$。
\end{frame}


\begin{frame}
  \frametitle{一般四维时空度规}
  \bmini{0.5}
  一般的四维时空是可以有内禀弯曲的，也就是说，无论如何取坐标系，度规 $g_{\mu\nu}$ 都不能处处和 $\eta_{\mu\nu}$ 相同:
  $$ ds^2= g_{\mu\nu}dx^\mu dx^\nu.$$
  \emini
  \bmini{0.45}
  \addfig{1.5}{curvedspace.jpg}
  \emini

  \skipline
  
  在宇宙的大多数地方，时空都比较接近 Minkowski 时空，可以把 $g_{\mu\nu}$ 写成
  \be
  g_{\mu\nu} = \eta_{\mu\nu} + h_{\mu\nu}.
  \ee
  这里的 $\lvert h_{\mu\nu}\rvert\ll 1$. 我们后面会专门对这种形式的度规进行研究，并发展一些非常有用的计算方法。

  \skipline

  在另外一些极端的地方，如黑洞、中子星等致密星体的周围。时空和 Minkowski 时空差距巨大。这种问题一般很难求解，我们将仅讨论一些最简单的情况。
  
\end{frame}


\secpage{联络}{$$\Gamma_{\lambda\mu\nu} = \frac{1}{2}\left(g_{\lambda\mu,\nu}+g_{\lambda\nu,\mu}-g_{\mu\nu,\lambda}\right).$$}

\begin{frame}
  我们先不假设度规的存在（因而也没有指标的升降等操作）。协变和逆变张量纯粹由数学定义给出，即它们满足张量的坐标变换性质
  {\blue  $$ \widetilde{T}_{\mu_1\mu_2\ldots}^{\ \  \nu_1\nu_2\ldots}= \frac{\partial x^{\lambda_1}}{\partial \tilde{x}^{\mu_1}} \frac{\partial x^{\lambda_2}}{\partial \tilde{x}^{\mu_2}}\ldots \frac{\partial \tilde{x}^{\nu_1}}{\partial x^{\rho_1}} \frac{\partial \tilde{x}^{\nu_2}}{\partial x^{\rho_2}} \ldots T_{\lambda_1\lambda_2\ldots}^{\ \ \rho_1\rho_2\ldots}.$$}
    即可。

    \skipline
    
    在此前提下，我们根本不讨论物理实体，而是把协变张量和逆变张量看成不同的数学对象。
\end{frame}


\begin{frame}
  \frametitle{联络的坐标变换性质}
  如果要求协变微商保持张量性，就可以推出，当从一个坐标系 $(x^0,x^1,x^2,x^3)$ 变换到另一个坐标系 $(\tilde{x}^0,\tilde{x}^1,\tilde{x}^2,\tilde{x}^3)$，联络的变换规则为：
  {\blue $$\tilde{\Gamma}^\lambda_{\ \mu\nu} = \frac{\partial \tilde{x}^\lambda}{\partial x^\alpha}\frac{\partial^2x^\alpha}{\partial \tilde{x}^\mu\partial \tilde{x}^\nu} + \Gamma^\alpha_{\ \beta\gamma}\frac{\partial \tilde{x}^\lambda}{\partial x^\alpha}\frac{\partial x^\beta}{\partial \tilde{x}^\mu}\frac{\partial x^\gamma}{\partial \tilde{x}^\nu} $$}
  由此可见，联络不是一个张量。

  \skipline

  如果有多种可能的联络，显然{\blue 任何两个联络的差是张量} (因为多余的那项 $\frac{\partial \tilde{x}^\lambda}{\partial x^\alpha}\frac{\partial^2x^\alpha}{\partial \tilde{x}^\mu\partial \tilde{x}^\nu}$ 抵消了)。反过来，{\blue 给联络 $\Gamma^\lambda_{\ \mu\nu}$ 加上任何张量 $B^\lambda_{\ \mu\nu}$ 都不会影响联络的坐标变换规则。}
\end{frame}



\begin{frame}
  \frametitle{协变微商的保张量性}
  我们再来回忆下协变微商公式
  {\blue
  \bea
  \left(T^{\mu_1\mu_2\ldots}_{\ \ \nu_1\nu_2\ldots}\right)_{;\lambda} &=& \left(T^{\mu_1\mu_2\ldots}_{\ \ \nu_1\nu_2\ldots}\right)_{,\lambda} \newl
  && + \Gamma^{\mu_1}_{\ \rho\lambda} T^{\rho\mu_2\ldots}_{\ \ \nu_1\nu_2\ldots} + \Gamma^{\mu_2}_{\ \rho\lambda} T^{\mu_1\rho\ldots}_{\ \ \nu_1\nu_2\ldots} + \ldots \newl
  && - \Gamma^\rho_{\nu_1\lambda}T^{\mu_1\mu_2\ldots}_{\ \ \rho\nu_2\ldots} - \Gamma^\rho_{\nu_2\lambda}T^{\mu_1\mu_2\ldots}_{\ \ \nu_1\rho\ldots}-\ldots
  \eea
  }
  显然，{\blue 给联络 $\Gamma^\lambda_{\ \mu\nu}$ 加上任何张量 $B^\lambda_{\ \mu\nu}$ 都不会破坏协变微商的保张量性。} 
\end{frame}


\begin{frame}
  现在我们引入度规

  \addfig{1.}{lundaowo.jpg}
\end{frame}


\begin{frame}
  \frametitle{协变微商的物理实体性}
  我们来看 $A^\mu_{\ ;\nu}$ 和 $A_{\mu;\nu}$ 之间的联系。

  \skiplines
  
  在曲面嵌入高维平直空间的假设下，矢量的协变微商具有明确的物理实体对应， $A^\mu_{\ ;\nu}$ 和 $A_{\mu;\nu}$ 对应同一物理实体，所以必然有 $A_{\mu;\nu} = g_{\mu\lambda}A^{\lambda}_{\ ;\nu}$。

  \skiplines

  如果我们进行的是抽象的不依赖于高维平直空间的讨论。那么除了我们恰好用了很具有迷惑性的符号之外，两者其实并无确定的关系。也就是说，如果把 $A_{\mu;\nu}$ 写成 $T_{\mu\nu}$，$A^\mu_{\ ;\nu}$ 写成 $S^\mu_{\ \nu}$，则 $T$ 和 $S$ 虽然都是张量（这由联络的坐标变换规则和协变微商的计算规则即可保证），但 $T$ 和 $S$ 未必对应同一物理量。

\end{frame}


\begin{frame}
  \frametitle{协变微商的物理实体性和度量的普适性之间的关系}
  显然，协变微商具有物理实体性（升降指标和求协变微商次序可以交换）和 度量的普适性(度规的协变微商处处为零)是同一件事情。

  \skiplines
  
  例如，如果度规的协变形式的协变微商为零，就有
  $$A_{\mu;\nu} = \left(g_{\mu\lambda}A^\lambda\right)_{;\nu} = g_{\mu\lambda;\nu}A^\lambda + g_{\mu\lambda}A^\lambda_{\ ;\nu}=g_{\mu\lambda}A^\lambda_{\ ;\nu}.$$
\end{frame}

\begin{frame}
  \frametitle{度量的普适性的等价条件}
  由于 $g^\mu_{\ \nu}$ 是常数矩阵，所以其普通微商为零：
  $$g^\mu_{\ \nu;\lambda} = \Gamma^\mu_{\ \alpha\lambda}g^\alpha_{\ \nu} - \Gamma^\alpha_{\ \nu\lambda}g^\mu_{\ \alpha} = \Gamma^\mu_{\ \nu\lambda} - \Gamma^\mu_{\ \nu\lambda} = 0.$$
  即度规的混合形式的协变微商自动为零，并不给出任何额外限制。
\end{frame}


\begin{frame}
  \frametitle{度量的普适性的等价条件(续)}
  度规的协变形式的协变微商为零给出：
  $$g_{\mu\nu;\lambda} = g_{\mu\nu,\lambda}-g_{\alpha\nu}\Gamma^\alpha_{\ \mu\lambda}-g_{\mu\alpha}\Gamma^\alpha_{\ \nu\lambda}=0.$$
  
  如果上述条件得到满足，我们并不需要额外考虑“度规的逆变形式的协变微商为零”这个条件，这是因为 $g^{\mu\nu} = g^\mu_{\ \lambda}g^\nu_{\ \rho} g_{\lambda\rho}$，而三个因子的协变微商均为零。
\end{frame}

\begin{frame}
  \frametitle{Christoffel-Levi-Civita联络}
    如果空间是 $n$ 维的，则 $g_{\mu\nu;\lambda}= 0$ 给出了 $\frac{n^2(n+1)}{2}$ 个限制方程，还不足以确定 $n^3$ 个联络。不过，{\blue 如果假设联络是对称的 $\Gamma^\lambda_{\ \mu\nu}=\Gamma^\lambda_{\ \nu\mu}$}，就可以通过直接化简 $g_{\mu\lambda;\nu}+g_{\nu\lambda;\mu}-g_{\mu\nu;\lambda}=0$ 得到
  {\blue  $$\Gamma_{\lambda\mu\nu}\equiv g_{\lambda\alpha}\Gamma^\alpha_{\ \mu\nu}=\frac{1}{2}\left(g_{\lambda\mu,\nu}+g_{\lambda\nu,\mu}-g_{\mu\nu,\lambda}\right);$$}
  它有个你记不住的名字，叫{\blue Christoffel-Levi-Civita联络，简称 Christoffel 联络}。

  \addfig{0.9}{levicivita.png}
\end{frame}


\begin{frame}
  \frametitle{挠还是不挠？}
  在一般的（不假设嵌入高维空间的图景）理论中，可以在联络里加入额外的一个反对称张量，称为{\blue 挠率张量}。这个挠率张量和曲线的挠率同样有“扭曲”的意思，但本质并不相同：前者描述的是由于空间坐标轴的代数不对称性造成的二阶小量修正，后者描述的是普通欧氏空间里的三阶小量修正。

  \addfig{2}{naonaonao.jpg}
  
  也许是基于我们的时空嵌入在一个高维平直空间里的大胆假设，或者纯粹为了减少理论的自由度，爱因斯坦在{\blue 广义相对论里假设不存在挠率张量，只使用 Cristoffel 联络}。
\end{frame}

\begin{frame}
  \frametitle{测地线}
  由于广义相对论使用 Cristoffel 联络，我们之前对测地线方程的推导完全可以照搬到四维时空来，结果同样是
  \tbox{$$\frac{d^2x^\lambda}{ds^2} + \Gamma^\lambda_{\ \mu\nu}\frac{dx^\mu}{ds}\frac{dx^\nu}{ds} = 0.$$}
  如果说广义相对论分为运动学和引力两个部分，那么上面的测地线方程就是运动学的核心方程。
\end{frame}


\begin{frame}
  在下一讲，我们将转入对四维时空的黎曼张量，为讨论广义相对论的第二核心方程——爱因斯坦引力方程做好准备。
\end{frame}

\ech
\end{document}


